
\section{Conclusion and Discussion}

We propose a Privacy-Preserving Interpretable Single-Domain Generalization framework for multi-label ocular disease recognition. The framework provides a unified solution to a set of progressive challenges: clinical multi-disease co-occurrence, cross-institutional domain shift, and the privacy sensitivity that further restricts multi-source collaboration. The core contributions are as follows: (1) the Adaptive Label-Feature Relevance Construction (ARC) module uses Transformer self-attention and cross-attention mechanisms to capture label-label relevance and label-feature relevance, and aligns dual-branch outputs through attention consistency loss, helping reduce false positives and false negatives in multi-label recognition; (2) the Dual-Branch Guided Domain-Invariant Feature Extraction (DFE) module simulates potential domain shift through style augmentation, uses CAM to provide disease-region prior guidance, and combines spatial/channel attention with a dual-branch mask-based contrastive learning mechanism to extract domain-invariant features, improving unseen-domain generalization performance while enhancing model interpretability; and (3) the Adaptive Gaussian Differential Privacy (AdaGDP) mechanism dynamically decays noise intensity to achieve a balance between privacy preservation and model performance under a moderate privacy budget (\(\epsilon \approx 10\)). Systematic experiments validate the effectiveness of the proposed method. Despite these advances, the framework still has limitations. Class imbalance may affect the stability of recognition for minority disease classes, which could be further alleviated by resampling or loss weighting in future work. In addition, current cross-domain generalization is mainly based on shared label spaces, and future work can explore generalization under open-category and more complex out-of-distribution scenarios. Furthermore, in terms of privacy preservation, future work can explore methods that retain performance under tighter privacy budgets, or combine the proposed mechanism with federated learning and similar techniques to adapt to more complex cross-institutional deployment scenarios.
